\begin{proofbox}
\textbf{\textcolor{CProof}{思路提示}}

利用等比数列通项公式将乘积化为含公比的幂形式，比较指数。

\medskip
\textbf{\textcolor{CProof}{正式推导}}
\begin{description}[style=nextline, leftmargin=2.8em, labelsep=0.5em, itemsep=0.55em]
    \item[\textbf{步骤一（设通项）：}] 设等比数列首项为 $a_1$，公比为 $q$，则对任意 $k\in\mathbb{N}^*$ 有
    \[
    a_k = a_1 q^{k-1}.
    \]
    \par\textbf{理由：} 等比数列通项公式。
    
    \item[\textbf{步骤二（写乘积）：}] 写出 $a_m\cdot a_n$ 和 $a_p\cdot a_q$ 的表达式：
    \[
    a_m\cdot a_n = (a_1 q^{m-1})(a_1 q^{n-1}) = a_1^2 q^{m+n-2},
    \]
    \[
    a_p\cdot a_q = (a_1 q^{p-1})(a_1 q^{q-1}) = a_1^2 q^{p+q-2}.
    \]
    \par\textbf{理由：} 代入通项并合并幂次。
    
    \item[\textbf{步骤三（比较指数）：}] 已知 $m+n = p+q$，因此 $m+n-2 = p+q-2$，故
    \[
    a_1^2 q^{m+n-2} = a_1^2 q^{p+q-2}.
    \]
    \par\textbf{理由：} 底数相同，指数相等则幂相等。
\end{description}

\medskip
\textbf{\textcolor{CProof}{结论回扣}}

由以上得 $a_m\cdot a_n = a_p\cdot a_q$，结论成立。
\end{proofbox}
